
\documentclass{article}
\usepackage{colortbl}
\usepackage{makecell}
\usepackage{multirow}
\usepackage{supertabular}

\begin{document}

\newcounter{utterance}

\centering \large Interaction Transcript for game `wordle\_withcritic', experiment `medium\_frequency\_words\_clue\_with\_critic', episode 5 with Qwen3.5{-}4B.
\vspace{24pt}

{ \footnotesize  \setcounter{utterance}{1}
\setlength{\tabcolsep}{0pt}
\begin{supertabular}{c@{$\;$}|p{.15\linewidth}@{}p{.15\linewidth}p{.15\linewidth}p{.15\linewidth}p{.15\linewidth}p{.15\linewidth}}
    \# & $\;$A & \multicolumn{4}{c}{Game Master} & $\;\:$B\\
    \hline

    \theutterance \stepcounter{utterance}  
    & & \multicolumn{4}{p{0.6\linewidth}}{
        \cellcolor[rgb]{0.9,0.9,0.9}{
            \makecell[{{p{\linewidth}}}]{
                \texttt{\tiny{[P1$\langle$GM]}}
                \texttt{You are a language wizard who likes to guess words by using the given rules.} \\
\\ 
\texttt{Welcome to Wordle! You have six attempts to guess the target word, a valid English word of five lowercase letters (a{-}z). Please use the tags "explanation:" and "guess:" to provide a concise explanation for each guess.} \\
\\ 
\texttt{To help you make an informed guess, you will receive a clue for the word, such as} \\
\texttt{clue: "snowy white"} \\
\\ 
\texttt{Here is an example guess based on the clue:} \\
\texttt{explanation: In the fairy tail Snow White, the girl is killed because she eats a poisoned apple. And the word apple has 5 letters.} \\
\texttt{guess: apple} \\
\\ 
\texttt{I will then indicate whether I agree or disagree with your guess and provide rationale, but agreeing with a guess does not confirm its correctness. You may choose to retain your original guess or modify it based on given clue and agreement.} \\
\\ 
\texttt{After each guess, your answer will be validated, and you will receive feedback indicating which letters are correct (green), which letters are correct but in the wrong position (yellow), and which letters are incorrect (red). This feedback can be useful in determining which letters to include or exclude in your next guess.} \\
\\ 
\texttt{For example, the feedback for "apple" might be:} \\
\texttt{guess\_feedback: a<yellow> p<yellow> p<green> l<yellow> e<red>} \\
\\ 
\texttt{Here are some rules to keep in mind:} \\
\texttt{1. Please do not generate any guess feedback or clue on your own} \\
\texttt{2. Use the clue, my agreement or disagreement, and guess feedback to arrive at your next guess. This should be reflected in your explanation as well.} \\
\\ 
\texttt{Do not generate any other text.} \\
\\ 
\texttt{Let's begin with your first guess.} \\
\\ 
\texttt{clue: tick off} \\
            }
        }
    }
    & \\ \\

    \theutterance \stepcounter{utterance}  
    & \multicolumn{4}{p{0.6\linewidth}}{
        \cellcolor[rgb]{0.95,0.95,0.95}{
            \makecell[{{p{\linewidth}}}]{
                \texttt{\tiny{[P2$\rightarrow$P1]}}
                \texttt{{[}DRAFT{]}} \\
\texttt{explanation: The clue "tick off" suggests a word related to marking something as done or completed. A common five{-}letter word fitting this description is "done up," but that's two words. Another possibility is "checked," which means to mark as done. Let's try "checked" as our first guess.} \\
\texttt{guess: checked} \\
            }
        }
    }
    & & \\ \\

    \theutterance \stepcounter{utterance}  
    & \multicolumn{4}{p{0.6\linewidth}}{
        \cellcolor[rgb]{0.95,0.95,0.95}{
            \makecell[{{p{\linewidth}}}]{
                \texttt{\tiny{[P2$\rightarrow$P1]}}
                \texttt{{[}SELF{-}FEEDBACK{]}} \\
\texttt{FORMAT: problem} \\
\texttt{RULES: problem} \\
\texttt{STRATEGY: problem} \\
\texttt{REASON: The answer uses incorrect tags ("explanation:" and "guess:" instead of "explanation:" and "guess:" as per the prompt's specific instruction to use those exact tags, though the tags themselves are correct, the explanation contains a phrase "Let's try 'checked' as our first guess" which violates the rule to not generate any guess feedback or clue on your own, and the strategy is poor because "checked" is not a valid five{-}letter word in standard English word lists for Wordle).} \\
\\ 
\texttt{Wait, looking closer at the prompt's instruction: "Please use the tags "explanation:" and "guess:"". The draft uses "explanation:" and "guess:". That part is correct.} \\
\texttt{The draft says: "guess: checked". "Checked" is 7 letters. Wordle requires 5 letters.} \\
\texttt{The draft explanation says: "Another possibility is 'checked,' which means to mark as done."} \\
\texttt{The draft explanation says: "Let's try 'checked' as our first guess."} \\
\\ 
\texttt{Let's re{-}evaluate based on the specific constraints.} \\
\texttt{1. FORMAT: The tags "explanation:" and "guess:" are used. No extra text outside the tags? The draft starts with "explanation:" and ends with "guess:". It seems to follow the format structure. However, the prompt says "Do not generate any other text." The draft does not have text outside the tags. So FORMAT might be ok.} \\
\texttt{2. RULES: The word "checked" has 7 letters. The game requires 5 letters. This is a violation of the game rules. Also, the explanation says "checked" is a possibility, but then guesses it. The explanation also says "done up" is two words. The main issue is the letter count.} \\
\texttt{3. STRATEGY: Guessing a 7{-}letter word is invalid, so strategy is moot, but guessing a word that doesn't fit the length constraint is a fundamental error.} \\
\\ 
\texttt{Let's refine the FORMAT check. The prompt says: "Please use the tags "explanation:" and "guess:" to provide a concise explanation for each guess." The draft does exactly that. It does not have extra text. So FORMAT is likely "ok".} \\
\texttt{RULES: "checked" is 7 letters. Wordle is 5 letters. This is a "problem".} \\
\texttt{STRATEGY: Since the guess is invalid, it's a "problem".} \\
\\ 
\texttt{Wait, looking at the example outputs in the prompt description:} \\
\texttt{"explanation: In the fairy tail Snow White, the girl is killed because she eats a poisoned apple. And the word apple has 5 letters."} \\
\texttt{"guess: apple"} \\
\\ 
\texttt{The draft:} \\
\texttt{"explanation: The clue "tick off" suggests a word related to marking something as done or completed. A common five{-}letter word fitting this description is "done} \\
            }
        }
    }
    & & \\ \\

    \theutterance \stepcounter{utterance}  
    & \multicolumn{4}{p{0.6\linewidth}}{
        \cellcolor[rgb]{0.9,0.9,0.9}{
            \makecell[{{p{\linewidth}}}]{
                \texttt{\tiny{[P1$\rangle$GM]}}
                \texttt{explanation: The clue "tick off" suggests a word related to marking something as done. A common five{-}letter word fitting this description is "marked".} \\
\texttt{guess: marked} \\
            }
        }
    }
    & & \\ \\

    \theutterance \stepcounter{utterance}  
    & & & \multicolumn{2}{p{0.3\linewidth}}{
        \cellcolor[rgb]{0.9,0.9,0.9}{
            \makecell[{{p{\linewidth}}}]{
                \texttt{\tiny{[GM$|$GM]}}
                \texttt{Error: The length of the guessed word is not 5.} \\
            }
        }
    }
    & & \\ \\

    \theutterance \stepcounter{utterance}  
    & & & \multicolumn{2}{p{0.3\linewidth}}{
        \cellcolor[rgb]{0.9,0.9,0.9}{
            \makecell[{{p{\linewidth}}}]{
                \texttt{\tiny{[GM$|$GM]}}
                \texttt{game\_result = ABORT} \\
            }
        }
    }
    & & \\ \\

\end{supertabular}
}

\end{document}
